% !TEX program = lualatex
% 自動生成: mar.report.latex — 手で編集した場合は再生成で上書きされる
\documentclass[paper=a4,fontsize=10.5pt,jafontsize=10.5pt]{jlreq}

\usepackage{luatexja}
\usepackage[haranoaji]{luatexja-preset}
\usepackage{amsmath,amssymb,amsthm,mathtools}
\usepackage{bm}
\usepackage{graphicx}
\usepackage{booktabs}
\usepackage{listings}
\usepackage{xcolor}
\usepackage[margin=25mm]{geometry}
\usepackage[hidelinks]{hyperref}
\usepackage{url}

\lstset{
  basicstyle=\ttfamily\footnotesize,
  breaklines=true,
  frame=single,
  showstringspaces=false,
  numbers=left,
  numberstyle=\tiny\color{gray},
  captionpos=b,
}

\theoremstyle{plain}
\newtheorem{theorem}{定理}[section]
\newtheorem{proposition}[theorem]{命題}
\newtheorem{lemma}[theorem]{補題}
\newtheorem{corollary}[theorem]{系}
\newtheorem{conjecture}[theorem]{予想}
\theoremstyle{definition}
\newtheorem{definition}[theorem]{定義}
\newtheorem{example}[theorem]{例}
\newtheorem{remark}[theorem]{注意}

\renewcommand{\proofname}{証明}

\title{飽和数は調和指数を超えうる: TxGraffiti 予想 4 の反例の完全分類}
\author{自動研究パイプライン \texttt{mar}（Claude Opus 5 + 検証器）}
\date{2026年07月26日}

\begin{document}
\maketitle

\begin{abstract}
自動予想生成系 TxGraffiti が提出した予想 (arXiv:2507.17780 Conjecture 4)、すなわち連結グラフ $G$ に対する $\mu^*(G) \le H(G)$ ($\mu^*$ は飽和数 = 最小極大マッチングの大きさ、$H$ は調和指数) は Bıyıkoğlu によって反証されている。本稿はこの反証を独立に追試したうえで、\emph{反例の完全な分類}を有限範囲で与える。すなわち $n \le 10$ の連結グラフ全 11,989,763 個と $n \le 20$ の木 全 1,345,823 個を走査し、$\mu^*(G) > H(G)$ となるグラフがちょうど 145 個であることを確定した。最小の反例は $n = 9$ の連結グラフ 8 個であり、木に限れば $n = 11$ のものが最小である。検証は p0002 と同じ\emph{証人つき}の設計をとる。各グラフに $|M| \le H(G)$ を満たす極大マッチング $M$ を添えれば $\mu^*(G) \le |M| \le H(G)$ が線形時間で確認でき、NP 困難な最小極大マッチングを検証器が解き直す必要がない。反例だと主張するグラフについてだけは、線グラフの極大独立集合を全列挙する独立実装で $\mu^*$ を厳密に計算し直している。各族の最大比についても、証人による上界と厳密再計算による下界を突き合わせて値を確定させた。
\end{abstract}

\noindent\textbf{キーワード:} graph theory、saturation number、maximal matching、harmonic index、TxGraffiti、refuted conjecture、classification、certificate


\section{はじめに}

グラフ $G$ のマッチング $M$ が\textbf{極大}であるとは、$M$ に含まれない
どの辺も $M$ の辺と端点を共有することをいう。極大マッチングの最小の大きさを
\textbf{飽和数} $\mu^*(G)$ と呼ぶ。$\mu^*(G)$ は最小辺支配集合の大きさに等しく、
その決定は 3-正則グラフに制限しても NP 困難である \cite{yannakakis}。
一方、\textbf{調和指数}は
\[ H(G) = \sum_{uv \in E(G)} \frac{2}{d(u) + d(v)} \]
で定義される次数だけから決まる量で、Fajtlowicz の Graffiti \cite{fajtlowicz} に
由来する化学グラフ理論の指標である。

TxGraffiti \cite{reverie} は両者を結ぶ次の予想を生成した。

\begin{conjecture}[TxGraffiti \cite{reverie}, Conjecture 4]
\label{conj:main}
連結グラフ $G$ に対し $\mu^*(G) \le H(G)$。
\end{conjecture}

予想 \ref{conj:main} は Bıyıkoğlu \cite{biyikoglu} によって反証された。
さらに Gupta \cite{gupta} は、木 $T$ について $\mu^*(T) < \tfrac{3}{2} H(T)$、
一般の連結グラフについて $H(G) < 4\mu^*(G)$ という両側の評価を与え、
最小の反例として 9 頂点のグラフと 11 頂点の細分星を挙げている。
しかし\emph{反例がどれだけあるのか}、すなわち小さい位数での反例の全体像は
報告されていない。本稿はそこを埋める。

\paragraph{本稿の貢献}
\begin{enumerate}
\item $n \le 10$ の連結グラフ全 11,989,763 個、および $n \le 20$ の木
  全 1,345,823 個について、予想 \ref{conj:main} の反例を\emph{すべて}決定した
  (定理 \ref{thm:main})。
\item 反例は合計 145 個で、その graph6 表記・次数列・$\mu^*$・$H$ を
  表 \ref{tab:ce} に列挙した。$n \le 8$ の連結グラフには反例が存在しない。
\item 比 $\mu^*(G)/H(G)$ の族内最大値を位数ごとに厳密に確定し、
  検証器がそれを上下から独立に閉じられる形にした (表 \ref{tab:fams})。
  観測された最大値は $\frac{33}{28}$ であり、Gupta の上界 $\tfrac{3}{2}$ とは
  まだ隔たりがある。
\item 検証を線形時間の証人検査に落とす設計を、この問題に対して具体化した
  (補題 \ref{lem:witness})。
\end{enumerate}

\section{証人つき片側検証}\label{sec:design}

反例でないこと、すなわち $\mu^*(G) \le H(G)$ を示すのに $\mu^*(G)$ を
求める必要はない。

\begin{lemma}\label{lem:witness}
$M$ が $G$ の極大マッチングで $|M| \le H(G)$ ならば $\mu^*(G) \le H(G)$ である。
\end{lemma}

\begin{proof}
$\mu^*(G)$ は極大マッチングの大きさの最小値だから $\mu^*(G) \le |M|$。
\end{proof}

補題 \ref{lem:witness} の仮定は、辺の重複がないこと・すべての辺が
$M$ の端点に接することの走査 ($M$ の極大性) と、次数から $H(G)$ を
足し合わせる計算だけで確認できる。どちらも辺数に関して線形である。
そこで本稿の証明書は、各グラフに対する $M$ そのものを証人とする。

逆向き、すなわち「これは反例である」という主張は片側評価では閉じない。
$\mu^*(G) > H(G)$ を示すには $\mu^*$ の\emph{下界}が要るからである。
そこで反例と主張するグラフに限り、検証器が $\mu^*$ を厳密に計算し直す。
反例は高々数十個なので、遅い独立実装で構わない。
実際、検証器は「$G$ の極大マッチング $\leftrightarrow$ 線グラフ $L(G)$ の
極大独立集合」という対応を使い、$L(G)$ の極大独立集合を
Bron--Kerbosch で全列挙して最小濃度を取る。探索器の分枝限定法とは
アルゴリズムが異なるので、共通のバグが両者をすり抜けることはない。

\subsection{探索の手順}

各グラフについて次を行う。

\begin{enumerate}
\item 次数の昇順、および頂点番号順の 2 通りで貪欲に極大マッチングを作り、
  小さいほうを $M$ とする。$|M| \le H(G)$ ならこれを証人として採用する。
\item そうでなければ最小極大マッチングを分枝限定法で厳密に求める。
  これで $\mu^*(G)$ が確定するので、$\mu^*(G) \le H(G)$ なら
  その最小極大マッチングを証人とし、そうでなければ反例として記録する。
\end{enumerate}

ただしこれだけでは族内の最大比 $\max \mu^*/H$ を取りこぼす。貪欲で
$|M| \le H$ が示せたグラフでも、真の $\mu^*$ が $|M|$ より小さいだけで
比が族内最大になる可能性があるからである。そこで探索器は現在の最大比を
保持し、証人の与える上界 $|M|/H$ がそれを上回るグラフについても厳密計算を
行う ($\mu^* \le |M|$ なので、上回らないグラフは族内最大になり得ない)。
最大比が $1$ に達した後はこの条件が「貪欲が $H$ を超えた」と一致するため、
追加コストは最大比が小さいうちの数十個に限られる。

貪欲で足りたグラフでは NP 困難な計算をまったく行わない。
貪欲が $H$ を超えたのは全 13,335,586 個中 11,249,412 個 (84.4\%) で、
最大比の追跡分を含めて厳密計算を行ったのは 11,295,265 個
(84.7\%) だった。

比較 $|M| \le H(G)$ は有理数演算で行うと遅い。位数 $n$ のグラフでは
$d(u)+d(v) \le 2n-2$ なので、$D_n = \mathrm{lcm}(2, 3, \dots, 2n-2)$ を共通分母に
とれば $H(G) \cdot D_n$ は整数であり、比較は整数演算だけで閉じる。
探索器はこの形で計算し、検証器は独立に \texttt{fractions.Fraction} で
厳密な有理数として計算し直す。

\subsection{証人の形式}

証人は族 (位数と種別の組) ごとに 1 つのバイナリファイルにまとめ、
gzip 圧縮して保存する。1 グラフあたり $n$ バイトを使い、
第 $i$ バイトは頂点 $i$ の相手の番号に 1 を足したもの
(マッチングに使われていなければ 0) とする。全バイトが 0 の記録は
「証人が存在しない = 反例である」という印である。
記録は元データの列挙順に並ぶので、検証器は元データを走査しながら
先頭から順に消費できる。各ファイルの SHA-256 は証明書 JSON に記録する。

\section{結果}

\begin{theorem}\label{thm:main}
$n \le 10$ の連結グラフ全 11,989,763 個と、$n \le 20$ の木
全 1,345,823 個のうち、$\mu^*(G) > H(G)$ を満たすものは
表 \ref{tab:ce} に挙げた 145 個ちょうどである。
とくに $n \le 8$ の連結グラフには反例が存在せず、
最小の反例は $n = 9$ の連結グラフ 8 個である。
木に限れば最小の反例は $n = 11$ にある。
\end{theorem}

\begin{proof}
表 \ref{tab:ce} に挙げた各グラフについては、$\mu^*$ を線グラフの
極大独立集合の全列挙で求め、$H$ を有理数演算で求めて
$\mu^* > H$ を直接確認した。それ以外のすべてのグラフについては、
$|M| \le H(G)$ を満たす極大マッチング $M$ を記録してあり、
補題 \ref{lem:witness} より反例ではない。
走査の網羅性は、各族で読み込んだ個数が OEIS A001349 (連結グラフ) および
A000055 (木) の公表値と一致すること (表 \ref{tab:fams}) による。
この公表値は探索器の出力ではなく、検証器が独自にもつ定数である。
これらの確認はすべて証明書 \texttt{p0003\_saturation\_harmonic.json} に対して
検証器が独立に実行した。
\end{proof}

\begin{table}[htbp]
\centering
\caption{走査した族。「厳密計算」は最小極大マッチングまで求めた個数
(貪欲で $|M| \le H$ を示せなかったグラフと、族内最大比を更新しうるグラフ)。
「最大比」は族内の $\mu^*(G)/H(G)$ の最大値で、厳密値である。}
\label{tab:fams}
\begin{tabular}{lrrrrr}
\toprule
種別 & $n$ & 個数 & 厳密計算 & 反例 & 最大比 \\
\midrule
連結グラフ & 2 & 1 & 1 & 0 & $1$ \\
連結グラフ & 3 & 2 & 1 & 0 & $\frac{3}{4}$ \\
連結グラフ & 4 & 6 & 5 & 0 & $1$ \\
連結グラフ & 5 & 21 & 8 & 0 & $\frac{70}{79}$ \\
連結グラフ & 6 & 112 & 101 & 0 & $1$ \\
連結グラフ & 7 & 853 & 207 & 0 & $1$ \\
連結グラフ & 8 & 11,117 & 10,001 & 0 & $1$ \\
連結グラフ & 9 & 261,080 & 36,064 & 8 & $\frac{10}{9}$ \\
連結グラフ & 10 & 11,716,571 & 10,967,170 & 6 & $\frac{220}{201}$ \\
木 & 11 & 235 & 85 & 1 & $\frac{21}{20}$ \\
木 & 12 & 551 & 160 & 1 & $\frac{420}{409}$ \\
木 & 13 & 1,301 & 370 & 2 & $\frac{12}{11}$ \\
木 & 14 & 3,159 & 804 & 2 & $\frac{72}{67}$ \\
木 & 15 & 7,741 & 1,935 & 4 & $\frac{9}{8}$ \\
木 & 16 & 19,320 & 4,750 & 7 & $\frac{315}{283}$ \\
木 & 17 & 48,629 & 11,497 & 9 & $\frac{15}{13}$ \\
木 & 18 & 123,867 & 29,760 & 19 & $\frac{1320}{1153}$ \\
木 & 19 & 317,955 & 64,175 & 37 & $\frac{33}{28}$ \\
木 & 20 & 823,065 & 168,171 & 49 & $\frac{198}{169}$ \\
\midrule
\multicolumn{2}{l}{合計} & 13,335,586 & 11,249,412 & 145 & \\
\bottomrule
\end{tabular}
\end{table}

\begin{table}[htbp]
\centering
\caption{予想 \ref{conj:main} の反例の完全な一覧
($n \le 10$ の連結グラフ、および $n \le 20$ の木)。}
\label{tab:ce}
\begin{tabular}{llrlrrr}
\toprule
graph6 & $n$ & $m$ & 次数列 & $\mu^*$ & $H$ & $\mu^*/H$ \\
\midrule
\texttt{H?`@?bw} & 9 & 9 & $(1,1,1,2,2,2,2,2,5)$ & 4 & $\frac{55}{14}$ & $\frac{56}{55}$ \\
\texttt{H?`@?b\{} & 9 & 10 & $(1,1,2,2,2,2,2,2,6)$ & 4 & $\frac{23}{6}$ & $\frac{24}{23}$ \\
\texttt{H?`@?b\}} & 9 & 11 & $(1,2,2,2,2,2,2,2,7)$ & 4 & $\frac{67}{18}$ & $\frac{72}{67}$ \\
\texttt{H?`@?b\textasciitilde{}} & 9 & 12 & $(2,2,2,2,2,2,2,2,8)$ & 4 & $\frac{18}{5}$ & $\frac{10}{9}$ \\
\texttt{H?`@`bz} & 9 & 13 & $(1,2,2,2,3,3,3,3,7)$ & 4 & $\frac{119}{30}$ & $\frac{120}{119}$ \\
\texttt{H?`@`b\}} & 9 & 13 & $(2,2,2,2,2,3,3,3,7)$ & 4 & $\frac{178}{45}$ & $\frac{90}{89}$ \\
\texttt{H?`@`b\textasciitilde{}} & 9 & 14 & $(2,2,2,2,3,3,3,3,8)$ & 4 & $\frac{637}{165}$ & $\frac{660}{637}$ \\
\texttt{HCOccf\textasciitilde{}} & 9 & 16 & $(2,2,2,2,4,4,4,4,8)$ & 4 & $\frac{119}{30}$ & $\frac{120}{119}$ \\
\texttt{I?AA@?O\textasciitilde{}\_} & 10 & 11 & $(1,1,1,2,2,2,2,2,2,7)$ & 4 & $\frac{47}{12}$ & $\frac{48}{47}$ \\
\texttt{I?AA@?O\textasciitilde{}o} & 10 & 12 & $(1,1,2,2,2,2,2,2,2,8)$ & 4 & $\frac{341}{90}$ & $\frac{360}{341}$ \\
\texttt{I?AA@?O\textasciitilde{}w} & 10 & 13 & $(1,2,2,2,2,2,2,2,2,9)$ & 4 & $\frac{201}{55}$ & $\frac{220}{201}$ \\
\texttt{I?AA@?W|w} & 10 & 13 & $(1,2,2,2,2,2,2,2,3,8)$ & 4 & $\frac{219}{55}$ & $\frac{220}{219}$ \\
\texttt{I?AA@?W\textasciitilde{}w} & 10 & 14 & $(2,2,2,2,2,2,2,2,3,9)$ & 4 & $\frac{647}{165}$ & $\frac{660}{647}$ \\
\texttt{I?AA@\_o\textasciitilde{}w} & 10 & 15 & $(1,2,2,2,2,3,3,3,3,9)$ & 4 & $\frac{216}{55}$ & $\frac{55}{54}$ \\
\texttt{J?AA@?OAFo?} & 11 & 10 & $(1,1,1,1,1,2,2,2,2,2,5)$ & 5 & $\frac{100}{21}$ & $\frac{21}{20}$ \\
\texttt{K??CA?\_C?O\textasciicircum{}\_} & 12 & 11 & $(1,1,1,1,1,1,2,2,2,2,2,6)$ & 5 & $\frac{409}{84}$ & $\frac{420}{409}$ \\
\texttt{L???C@?G?\_@?\textasciitilde{}\_} & 13 & 12 & $(1,1,1,1,1,1,1,2,2,2,2,2,7)$ & 5 & $\frac{89}{18}$ & $\frac{90}{89}$ \\
\texttt{L??CA?\_C?O?\_\textasciitilde{}?} & 13 & 12 & $(1,1,1,1,1,1,2,2,2,2,2,2,6)$ & 6 & $\frac{11}{2}$ & $\frac{12}{11}$ \\
\texttt{M???C@?G?\_@?@?\textasciitilde{}\_?} & 14 & 13 & $(1,1,1,1,1,1,1,2,2,2,2,2,2,7)$ & 6 & $\frac{67}{12}$ & $\frac{72}{67}$ \\
\texttt{M???C@?G?\_@?@\_\}?\_} & 14 & 13 & $(1,1,1,1,1,1,1,2,2,2,2,2,3,6)$ & 6 & $\frac{209}{36}$ & $\frac{216}{209}$ \\
\texttt{N????A?O@?A?A?@?\textasciicircum{}w?} & 15 & 14 & $(1,1,1,1,1,1,1,1,2,2,2,2,2,2,8)$ & 6 & $\frac{254}{45}$ & $\frac{135}{127}$ \\
\texttt{N????A?O@?A?A?@\_\textasciicircum{}GG} & 15 & 14 & $(1,1,1,1,1,1,1,1,2,2,2,2,2,3,7)$ & 6 & $\frac{1061}{180}$ & $\frac{1080}{1061}$ \\
\texttt{N????A?O@?A?A?@o\textasciicircum{}?G} & 15 & 14 & $(1,1,1,1,1,1,1,1,2,2,2,2,2,4,6)$ & 6 & $\frac{359}{60}$ & $\frac{360}{359}$ \\
\texttt{N???C@?G?\_@?@??\_\textasciicircum{}o?} & 15 & 14 & $(1,1,1,1,1,1,1,2,2,2,2,2,2,2,7)$ & 7 & $\frac{56}{9}$ & $\frac{9}{8}$ \\
\texttt{O??????\_A?C?C?A??\_F\textasciitilde{}?} & 16 & 15 & $(1,1,1,1,1,1,1,1,1,2,2,2,2,2,2,9)$ & 6 & $\frac{313}{55}$ & $\frac{330}{313}$ \\
\texttt{O??????\_A?C?C?A??oFr@} & 16 & 15 & $(1,1,1,1,1,1,1,1,1,2,2,2,2,2,3,8)$ & 6 & $\frac{590}{99}$ & $\frac{297}{295}$ \\
\texttt{O????A?O@?A?A?@??OF\}?} & 16 & 15 & $(1,1,1,1,1,1,1,1,2,2,2,2,2,2,2,8)$ & 7 & $\frac{283}{45}$ & $\frac{315}{283}$ \\
\texttt{O????A?O@?A?A?@??WFw@} & 16 & 15 & $(1,1,1,1,1,1,1,1,2,2,2,2,2,2,3,7)$ & 7 & $\frac{98}{15}$ & $\frac{15}{14}$ \\
\texttt{O????A?O@?A?A?@??WF\{?} & 16 & 15 & $(1,1,1,1,1,1,1,2,2,2,2,2,2,2,2,7)$ & 7 & $\frac{121}{18}$ & $\frac{126}{121}$ \\
\texttt{O????A?O@?A?A?@\_?gFoA} & 16 & 15 & $(1,1,1,1,1,1,1,2,2,2,2,2,2,2,3,6)$ & 7 & $\frac{1237}{180}$ & $\frac{1260}{1237}$ \\
\texttt{O????A?O@?A?A?@\_?gFw?} & 16 & 15 & $(1,1,1,1,1,1,1,2,2,2,2,2,2,2,3,6)$ & 7 & $\frac{1249}{180}$ & $\frac{1260}{1249}$ \\
\texttt{P???????C?G?G?C?@??G?\textasciitilde{}\{?} & 17 & 16 & $(1,1,1,1,1,1,1,1,1,1,2,2,2,2,2,2,10)$ & 6 & $\frac{63}{11}$ & $\frac{22}{21}$ \\
\texttt{P??????\_A?C?C?A??\_?C?\textasciitilde{}w?} & 17 & 16 & $(1,1,1,1,1,1,1,1,1,2,2,2,2,2,2,2,9)$ & 7 & $\frac{1046}{165}$ & $\frac{1155}{1046}$ \\
\texttt{P??????\_A?C?C?A??\_?E?\textasciitilde{}GC} & 17 & 16 & $(1,1,1,1,1,1,1,1,1,2,2,2,2,2,2,3,8)$ & 7 & $\frac{3269}{495}$ & $\frac{495}{467}$ \\
\texttt{P??????\_A?C?C?A??\_?E?\textasciitilde{}g?} & 17 & 16 & $(1,1,1,1,1,1,1,1,2,2,2,2,2,2,2,2,8)$ & 7 & $\frac{611}{90}$ & $\frac{630}{611}$ \\
\texttt{P??????\_A?C?C?A??\_?F?\textasciitilde{}?C} & 17 & 16 & $(1,1,1,1,1,1,1,1,1,2,2,2,2,2,2,4,7)$ & 7 & $\frac{1108}{165}$ & $\frac{1155}{1108}$ \\
\texttt{P??????\_A?C?C?A??\_?F?\textasciitilde{}\_?} & 17 & 16 & $(1,1,1,1,1,1,1,1,2,2,2,2,2,2,2,3,7)$ & 7 & $\frac{313}{45}$ & $\frac{315}{313}$ \\
\texttt{P??????\_A?C?C?A??o?B?\}?K} & 17 & 16 & $(1,1,1,1,1,1,1,1,1,2,2,2,2,2,3,3,7)$ & 7 & $\frac{308}{45}$ & $\frac{45}{44}$ \\
\texttt{P??????\_A?C?C?A??o?I?\}GG} & 17 & 16 & $(1,1,1,1,1,1,1,1,2,2,2,2,2,2,2,3,7)$ & 7 & $\frac{1253}{180}$ & $\frac{180}{179}$ \\
\texttt{P????A?O@?A?A?@??O?A?\textasciitilde{}o?} & 17 & 16 & $(1,1,1,1,1,1,1,1,2,2,2,2,2,2,2,2,8)$ & 8 & $\frac{104}{15}$ & $\frac{15}{13}$ \\
\texttt{Q?????????O?O?G?A??O?@?B\textasciitilde{}w?} & 18 & 17 & $(1,1,1,1,1,1,1,1,1,1,1,2,2,2,2,2,2,11)$ & 6 & $\frac{449}{78}$ & $\frac{468}{449}$ \\
\texttt{Q???????C?G?G?C?@??G??\_B\textasciitilde{}o?} & 18 & 17 & $(1,1,1,1,1,1,1,1,1,1,2,2,2,2,2,2,2,10)$ & 7 & $\frac{421}{66}$ & $\frac{462}{421}$ \\
\texttt{Q???????C?G?G?C?@??G??oB\{oG} & 18 & 17 & $(1,1,1,1,1,1,1,1,1,1,2,2,2,2,2,2,3,9)$ & 7 & $\frac{2197}{330}$ & $\frac{2310}{2197}$ \\
\texttt{Q???????C?G?G?C?@??G??oB\}o?} & 18 & 17 & $(1,1,1,1,1,1,1,1,1,2,2,2,2,2,2,2,2,9)$ & 7 & $\frac{2257}{330}$ & $\frac{2310}{2257}$ \\
\texttt{Q???????C?G?G?C?@??G??wB\{OG} & 18 & 17 & $(1,1,1,1,1,1,1,1,1,1,2,2,2,2,2,2,4,8)$ & 7 & $\frac{611}{90}$ & $\frac{630}{611}$ \\
\texttt{Q???????C?G?G?C?@??G??\{B\{?G} & 18 & 17 & $(1,1,1,1,1,1,1,1,1,1,2,2,2,2,2,2,5,7)$ & 7 & $\frac{41}{6}$ & $\frac{42}{41}$ \\
\texttt{Q???????C?G?G?C?@??K??WBwOW} & 18 & 17 & $(1,1,1,1,1,1,1,1,1,1,2,2,2,2,2,3,3,8)$ & 7 & $\frac{685}{99}$ & $\frac{693}{685}$ \\
\texttt{Q??????\_A?C?C?A??\_?C??OB\textasciitilde{}\_?} & 18 & 17 & $(1,1,1,1,1,1,1,1,1,2,2,2,2,2,2,2,2,9)$ & 8 & $\frac{1153}{165}$ & $\frac{1320}{1153}$ \\
\texttt{Q??????\_A?C?C?A??\_?C??WB\}?G} & 18 & 17 & $(1,1,1,1,1,1,1,1,1,2,2,2,2,2,2,2,3,8)$ & 8 & $\frac{1196}{165}$ & $\frac{330}{299}$ \\
\texttt{Q??????\_A?C?C?A??\_?C??WB\textasciitilde{}??} & 18 & 17 & $(1,1,1,1,1,1,1,1,2,2,2,2,2,2,2,2,2,8)$ & 8 & $\frac{223}{30}$ & $\frac{240}{223}$ \\
\texttt{Q??????\_A?C?C?A??\_?E??gB\{?O} & 18 & 17 & $(1,1,1,1,1,1,1,1,2,2,2,2,2,2,2,2,3,7)$ & 8 & $\frac{38}{5}$ & $\frac{20}{19}$ \\
\texttt{Q??????\_A?C?C?A??\_?E??gB|??} & 18 & 17 & $(1,1,1,1,1,1,1,2,2,2,2,2,2,2,2,2,2,7)$ & 8 & $\frac{139}{18}$ & $\frac{144}{139}$ \\
\texttt{Q??????\_A?C?C?A??\_?E??gB\}??} & 18 & 17 & $(1,1,1,1,1,1,1,1,2,2,2,2,2,2,2,2,3,7)$ & 8 & $\frac{23}{3}$ & $\frac{24}{23}$ \\
\texttt{Q????A?O@?A?A?@??O?A?w??\textasciicircum{}?G} & 18 & 17 & $(1,1,1,1,1,1,1,1,2,2,2,2,2,2,2,2,4,6)$ & 8 & $\frac{467}{60}$ & $\frac{480}{467}$ \\
\texttt{Q????A?O@?A?A?@??O?A?w??\textasciicircum{}\_?} & 18 & 17 & $(1,1,1,1,1,1,1,2,2,2,2,2,2,2,2,2,3,6)$ & 8 & $\frac{118}{15}$ & $\frac{60}{59}$ \\
\texttt{Q????A?O@?A?A?@??O?A?w?A\textasciicircum{}??} & 18 & 17 & $(1,1,1,1,1,1,1,1,2,2,2,2,2,2,2,3,3,6)$ & 8 & $\frac{1399}{180}$ & $\frac{1440}{1399}$ \\
\texttt{Q????A?O@?A?A?@??O?A?\{??N?G} & 18 & 17 & $(1,1,1,1,1,1,1,1,2,2,2,2,2,2,2,2,5,5)$ & 8 & $\frac{821}{105}$ & $\frac{840}{821}$ \\
\texttt{Q????A?O@?A?A?@??O?A?\{??N\_?} & 18 & 17 & $(1,1,1,1,1,1,1,2,2,2,2,2,2,2,2,2,4,5)$ & 8 & $\frac{111}{14}$ & $\frac{112}{111}$ \\
\texttt{Q????A?O@?A?A?@??O?A?\{?AN??} & 18 & 17 & $(1,1,1,1,1,1,1,1,2,2,2,2,2,2,2,3,4,5)$ & 8 & $\frac{659}{84}$ & $\frac{672}{659}$ \\
\texttt{RhCS?C@?\_?\_O?@?\_??GC???CA????G} & 19 & 18 & $(1,1,1,1,1,1,1,1,2,2,2,2,2,2,2,2,2,3,7)$ & 8 & $\frac{358}{45}$ & $\frac{180}{179}$ \\
\texttt{RhC\_K?@?GA?@?O??\_G??@?O???GA??} & 19 & 18 & $(1,1,1,1,1,1,1,1,2,2,2,2,2,2,2,2,2,3,7)$ & 8 & $\frac{159}{20}$ & $\frac{160}{159}$ \\
\texttt{RhE?GCC?GG?@@???\_\_??@@????GG??} & 19 & 18 & $(1,1,1,1,1,1,1,1,2,2,2,2,2,2,2,2,2,2,8)$ & 8 & $\frac{701}{90}$ & $\frac{720}{701}$ \\
\texttt{RhH?I?@O??g??@O???K????C??K???} & 19 & 18 & $(1,1,1,1,1,1,1,1,1,2,2,2,2,2,2,2,3,3,7)$ & 8 & $\frac{356}{45}$ & $\frac{90}{89}$ \\
\texttt{RhH?K?@\_??o??@\_???K????E?????G} & 19 & 18 & $(1,1,1,1,1,1,1,1,1,2,2,2,2,2,2,2,3,3,7)$ & 8 & $\frac{119}{15}$ & $\frac{120}{119}$ \\
\texttt{RhI?GCC?GG?@@???\_\_??@@????K???} & 19 & 18 & $(1,1,1,1,1,1,1,1,1,2,2,2,2,2,2,2,3,3,7)$ & 8 & $\frac{47}{6}$ & $\frac{48}{47}$ \\
\texttt{RhI?K?@\_??o??@\_???K????E?????G} & 19 & 18 & $(1,1,1,1,1,1,1,1,1,2,2,2,2,2,2,2,2,3,8)$ & 8 & $\frac{23}{3}$ & $\frac{24}{23}$ \\
\texttt{RhOI?D??I??@O???g???@\_????G??G} & 19 & 18 & $(1,1,1,1,1,1,1,1,2,2,2,2,2,2,2,2,2,2,8)$ & 8 & $\frac{119}{15}$ & $\frac{120}{119}$ \\
\texttt{RhOI?D??I??@O???g?C????C??K???} & 19 & 18 & $(1,1,1,1,1,1,1,1,1,2,2,2,2,2,2,2,2,3,8)$ & 8 & $\frac{3797}{495}$ & $\frac{3960}{3797}$ \\
\texttt{RhOI?D??I??@O???o???@??C??c???} & 19 & 18 & $(1,1,1,1,1,1,1,1,1,2,2,2,2,2,2,2,3,3,7)$ & 8 & $\frac{119}{15}$ & $\frac{120}{119}$ \\
\texttt{RhOI?D??I??@O???o???@??E??C???} & 19 & 18 & $(1,1,1,1,1,1,1,1,1,2,2,2,2,2,2,2,2,4,7)$ & 8 & $\frac{2579}{330}$ & $\frac{2640}{2579}$ \\
\texttt{RhOI?D??I??@O?\_???G?@??O??K???} & 19 & 18 & $(1,1,1,1,1,1,1,1,1,2,2,2,2,2,2,2,3,3,7)$ & 8 & $\frac{283}{36}$ & $\frac{288}{283}$ \\
\texttt{RhOI?D??I??@\_???\_?G?C??C?AC???} & 19 & 18 & $(1,1,1,1,1,1,1,1,1,2,2,2,2,2,2,2,3,4,6)$ & 8 & $\frac{10069}{1260}$ & $\frac{10080}{10069}$ \\
\texttt{RhOI?D??I?O??@??\_?\_?@?@???K???} & 19 & 18 & $(1,1,1,1,1,1,1,1,1,2,2,2,2,2,2,2,3,4,6)$ & 8 & $\frac{1003}{126}$ & $\frac{1008}{1003}$ \\
\texttt{RhOI?D??K??@\_???o???@\_????K???} & 19 & 18 & $(1,1,1,1,1,1,1,1,1,2,2,2,2,2,2,2,2,5,6)$ & 8 & $\frac{1835}{231}$ & $\frac{1848}{1835}$ \\
\texttt{RhOI?D?\_??o??@\_???K????E?????G} & 19 & 18 & $(1,1,1,1,1,1,1,1,1,2,2,2,2,2,2,2,2,5,6)$ & 8 & $\frac{7351}{924}$ & $\frac{7392}{7351}$ \\
\texttt{RhOI?E??K??@\_???o???@\_????K???} & 19 & 18 & $(1,1,1,1,1,1,1,1,1,2,2,2,2,2,2,2,2,4,7)$ & 8 & $\frac{3119}{396}$ & $\frac{3168}{3119}$ \\
\texttt{RhOIC?@\_??o??@\_???K????E?????G} & 19 & 18 & $(1,1,1,1,1,1,1,1,1,2,2,2,2,2,2,2,2,4,7)$ & 8 & $\frac{1306}{165}$ & $\frac{660}{653}$ \\
\texttt{RhOK?CA\_??o??@\_???K????E?????G} & 19 & 18 & $(1,1,1,1,1,1,1,1,1,2,2,2,2,2,2,2,3,3,7)$ & 8 & $\frac{359}{45}$ & $\frac{360}{359}$ \\
\texttt{RhOK?E??K??@\_???o???@\_????K???} & 19 & 18 & $(1,1,1,1,1,1,1,1,1,2,2,2,2,2,2,2,2,3,8)$ & 8 & $\frac{766}{99}$ & $\frac{396}{383}$ \\
\texttt{RhQ?K?@\_??o??@\_???K????E?????G} & 19 & 18 & $(1,1,1,1,1,1,1,1,1,2,2,2,2,2,2,2,2,3,8)$ & 8 & $\frac{2579}{330}$ & $\frac{2640}{2579}$ \\
\texttt{Rh\_GK?@\_??o??@\_???K????E?????G} & 19 & 18 & $(1,1,1,1,1,1,1,1,2,2,2,2,2,2,2,2,2,2,8)$ & 8 & $\frac{119}{15}$ & $\frac{120}{119}$ \\
\texttt{Rh\_GS?@\_??o??@\_???K????E?????G} & 19 & 18 & $(1,1,1,1,1,1,1,1,1,2,2,2,2,2,2,2,2,3,8)$ & 8 & $\frac{2557}{330}$ & $\frac{2640}{2557}$ \\
\texttt{Rh\_K?E??K??@\_???o???@\_????K???} & 19 & 18 & $(1,1,1,1,1,1,1,1,1,2,2,2,2,2,2,2,2,2,9)$ & 8 & $\frac{2471}{330}$ & $\frac{2640}{2471}$ \\
\texttt{Rh\_K?E??K??@\_???o???@\_?A??C???} & 19 & 18 & $(1,1,1,1,1,1,1,1,1,1,2,2,2,2,2,2,2,2,10)$ & 7 & $\frac{227}{33}$ & $\frac{231}{227}$ \\
\texttt{RiPAC?@\_??o??@\_???K????E?????G} & 19 & 18 & $(1,1,1,1,1,1,1,1,1,1,1,2,2,2,2,2,2,6,7)$ & 7 & $\frac{1888}{273}$ & $\frac{1911}{1888}$ \\
\texttt{RiPC?E??K??@\_???o???@\_????K???} & 19 & 18 & $(1,1,1,1,1,1,1,1,1,1,1,2,2,2,2,2,2,5,8)$ & 7 & $\frac{4042}{585}$ & $\frac{4095}{4042}$ \\
\texttt{RiQ?K?@\_??o??@\_???K????E?????G} & 19 & 18 & $(1,1,1,1,1,1,1,1,1,1,2,2,2,2,2,2,2,4,8)$ & 8 & $\frac{223}{30}$ & $\frac{240}{223}$ \\
\texttt{RiQ?K?@\_??o??@\_???K????E??C???} & 19 & 18 & $(1,1,1,1,1,1,1,1,1,1,1,2,2,2,2,2,2,4,9)$ & 7 & $\frac{4894}{715}$ & $\frac{5005}{4894}$ \\
\texttt{Ri\_GS?@\_??o??@\_???K????E?????G} & 19 & 18 & $(1,1,1,1,1,1,1,1,1,1,2,2,2,2,2,2,3,3,8)$ & 8 & $\frac{416}{55}$ & $\frac{55}{52}$ \\
\texttt{Ri\_GS?@\_??o??@\_???K????E??C???} & 19 & 18 & $(1,1,1,1,1,1,1,1,1,1,1,2,2,2,2,2,3,3,9)$ & 7 & $\frac{1151}{165}$ & $\frac{1155}{1151}$ \\
\texttt{Ri\_K?E??K??@\_???o???@\_????K???} & 19 & 18 & $(1,1,1,1,1,1,1,1,1,1,2,2,2,2,2,2,2,3,9)$ & 8 & $\frac{2411}{330}$ & $\frac{2640}{2411}$ \\
\texttt{Ri\_K?E??K??@\_???o???@\_?A??C???} & 19 & 18 & $(1,1,1,1,1,1,1,1,1,1,1,2,2,2,2,2,2,3,10)$ & 7 & $\frac{958}{143}$ & $\frac{1001}{958}$ \\
\texttt{RkE?K?@\_??o??@\_???K????E?????G} & 19 & 18 & $(1,1,1,1,1,1,1,1,1,2,2,2,2,2,2,2,2,2,9)$ & 9 & $\frac{84}{11}$ & $\frac{33}{28}$ \\
\texttt{RkE?K?@\_??o??@\_???K????E??C???} & 19 & 18 & $(1,1,1,1,1,1,1,1,1,1,2,2,2,2,2,2,2,2,10)$ & 8 & $\frac{232}{33}$ & $\frac{33}{29}$ \\
\texttt{RkE?K?@\_??o??@\_???K??\_?A??C???} & 19 & 18 & $(1,1,1,1,1,1,1,1,1,1,1,2,2,2,2,2,2,2,11)$ & 7 & $\frac{250}{39}$ & $\frac{273}{250}$ \\
\texttt{RkE?K?@\_??o??@\_?O?C??\_?A??C???} & 19 & 18 & $(1,1,1,1,1,1,1,1,1,1,1,1,2,2,2,2,2,2,12)$ & 6 & $\frac{526}{91}$ & $\frac{273}{263}$ \\
\texttt{ShC\_H?@G??o??@??\_?\_?@?@???G?G???C} & 20 & 19 & $(1,1,1,1,1,1,1,1,2,2,2,2,2,2,2,2,2,2,5,5)$ & 9 & $\frac{365}{42}$ & $\frac{378}{365}$ \\
\texttt{ShC\_H?@\_??\_@?C??\_A??@?C???G?\_???C} & 20 & 19 & $(1,1,1,1,1,1,1,1,2,2,2,2,2,2,2,2,2,2,4,6)$ & 9 & $\frac{26}{3}$ & $\frac{27}{26}$ \\
\texttt{ShC\_K?@?GA?@?O??\_G??@?O???GA????C} & 20 & 19 & $(1,1,1,1,1,1,1,1,2,2,2,2,2,2,2,2,2,2,3,7)$ & 9 & $\frac{773}{90}$ & $\frac{810}{773}$ \\
\texttt{ShE?GCC?GG?@@???\_\_??@@????GG????C} & 20 & 19 & $(1,1,1,1,1,1,1,1,2,2,2,2,2,2,2,2,2,2,2,8)$ & 9 & $\frac{253}{30}$ & $\frac{270}{253}$ \\
\texttt{ShE?GCC?GG?@@???\_\_??@@????GG??G??} & 20 & 19 & $(1,1,1,1,1,1,1,1,1,2,2,2,2,2,2,2,2,2,2,9)$ & 8 & $\frac{2587}{330}$ & $\frac{2640}{2587}$ \\
\texttt{ShGc?C@?\_?\_O?@?\_??GC???CA????K???} & 20 & 19 & $(1,1,1,1,1,1,1,1,1,1,2,2,2,2,2,2,2,3,4,7)$ & 8 & $\frac{239}{30}$ & $\frac{240}{239}$ \\
\texttt{ShGc?E??K??@\_???o???@\_????K?????C} & 20 & 19 & $(1,1,1,1,1,1,1,1,1,1,2,2,2,2,2,2,2,2,4,8)$ & 8 & $\frac{39}{5}$ & $\frac{40}{39}$ \\
\texttt{ShH?I?@O??g??@O???I??\_????G??K???} & 20 & 19 & $(1,1,1,1,1,1,1,1,1,1,2,2,2,2,2,2,2,3,3,8)$ & 8 & $\frac{3953}{495}$ & $\frac{3960}{3953}$ \\
\texttt{ShI?GCC?GG?@@???\_\_??@@????GG?C???} & 20 & 19 & $(1,1,1,1,1,1,1,1,1,1,2,2,2,2,2,2,2,3,3,8)$ & 8 & $\frac{1565}{198}$ & $\frac{1584}{1565}$ \\
\texttt{ShI?GI??K??@\_???o???@\_????K?????C} & 20 & 19 & $(1,1,1,1,1,1,1,1,1,1,2,2,2,2,2,2,2,3,3,8)$ & 8 & $\frac{439}{55}$ & $\frac{440}{439}$ \\
\texttt{ShI?K?@\_??o??@\_???K????E?????K???} & 20 & 19 & $(1,1,1,1,1,1,1,1,1,1,2,2,2,2,2,2,2,2,3,9)$ & 8 & $\frac{1274}{165}$ & $\frac{660}{637}$ \\
\texttt{ShOI?D??I??@O???g???@O?A?????G??C} & 20 & 19 & $(1,1,1,1,1,1,1,1,1,2,2,2,2,2,2,2,2,2,2,9)$ & 8 & $\frac{1318}{165}$ & $\frac{660}{659}$ \\
\texttt{ShOI?D??I??@O???g???@\_????G??K???} & 20 & 19 & $(1,1,1,1,1,1,1,1,1,2,2,2,2,2,2,2,2,2,3,8)$ & 9 & $\frac{1372}{165}$ & $\frac{1485}{1372}$ \\
\texttt{ShOI?D??I??@O???g?A??\_????G??K???} & 20 & 19 & $(1,1,1,1,1,1,1,1,1,1,2,2,2,2,2,2,2,2,3,9)$ & 8 & $\frac{2549}{330}$ & $\frac{2640}{2549}$ \\
\texttt{ShOI?D??I??@O???g?C????C??K??C???} & 20 & 19 & $(1,1,1,1,1,1,1,1,1,1,2,2,2,2,2,2,2,2,4,8)$ & 8 & $\frac{71}{9}$ & $\frac{72}{71}$ \\
\texttt{ShOI?D??I??@O???o???@??C??\_??K???} & 20 & 19 & $(1,1,1,1,1,1,1,1,1,2,2,2,2,2,2,2,2,3,3,7)$ & 9 & $\frac{17}{2}$ & $\frac{18}{17}$ \\
\texttt{ShOI?D??I??@O???o???@??E??C??C???} & 20 & 19 & $(1,1,1,1,1,1,1,1,1,1,2,2,2,2,2,2,2,2,5,7)$ & 8 & $\frac{167}{21}$ & $\frac{168}{167}$ \\
\texttt{ShOI?D??I??@O?O?O???@??C??\_??K???} & 20 & 19 & $(1,1,1,1,1,1,1,1,1,1,2,2,2,2,2,2,2,3,3,8)$ & 8 & $\frac{7847}{990}$ & $\frac{7920}{7847}$ \\
\texttt{ShOI?D??I??@\_???\_?G?C??C?A???K???} & 20 & 19 & $(1,1,1,1,1,1,1,1,1,2,2,2,2,2,2,2,2,3,4,6)$ & 9 & $\frac{2165}{252}$ & $\frac{2268}{2165}$ \\
\texttt{ShOI?D??K??@?@?A??G?O??C?G???K???} & 20 & 19 & $(1,1,1,1,1,1,1,1,1,2,2,2,2,2,2,2,2,3,5,5)$ & 9 & $\frac{181}{21}$ & $\frac{189}{181}$ \\
\texttt{ShOI?D??K??@\_???o???@\_????K?????C} & 20 & 19 & $(1,1,1,1,1,1,1,1,1,2,2,2,2,2,2,2,2,2,5,6)$ & 9 & $\frac{2641}{308}$ & $\frac{2772}{2641}$ \\
\texttt{ShOI?E??K??@\_???o???@\_????K?????C} & 20 & 19 & $(1,1,1,1,1,1,1,1,1,2,2,2,2,2,2,2,2,2,4,7)$ & 9 & $\frac{281}{33}$ & $\frac{297}{281}$ \\
\texttt{ShOI?E??K??@\_???o???@\_????K??C???} & 20 & 19 & $(1,1,1,1,1,1,1,1,1,1,2,2,2,2,2,2,2,2,4,8)$ & 8 & $\frac{143}{18}$ & $\frac{144}{143}$ \\
\texttt{ShOIC?@\_??o??@\_???K????E?????K???} & 20 & 19 & $(1,1,1,1,1,1,1,1,1,1,2,2,2,2,2,2,2,2,4,8)$ & 8 & $\frac{719}{90}$ & $\frac{720}{719}$ \\
\texttt{ShOK?E??K??@\_???o???@\_????K?????C} & 20 & 19 & $(1,1,1,1,1,1,1,1,1,2,2,2,2,2,2,2,2,2,3,8)$ & 9 & $\frac{461}{55}$ & $\frac{495}{461}$ \\
\texttt{ShOK?E??K??@\_???o???@\_????K??C???} & 20 & 19 & $(1,1,1,1,1,1,1,1,1,1,2,2,2,2,2,2,2,2,3,9)$ & 8 & $\frac{857}{110}$ & $\frac{880}{857}$ \\
\texttt{ShQ?K?@\_??o??@\_???K????E?????K???} & 20 & 19 & $(1,1,1,1,1,1,1,1,1,1,2,2,2,2,2,2,2,2,3,9)$ & 8 & $\frac{433}{55}$ & $\frac{440}{433}$ \\
\texttt{Sh\_GK?@\_??o??@\_???K????E?????K???} & 20 & 19 & $(1,1,1,1,1,1,1,1,1,2,2,2,2,2,2,2,2,2,2,9)$ & 8 & $\frac{1318}{165}$ & $\frac{660}{659}$ \\
\texttt{Sh\_GOQ??K??@\_???o???@\_????K?????C} & 20 & 19 & $(1,1,1,1,1,1,1,1,1,1,2,2,2,2,2,2,2,2,4,8)$ & 8 & $\frac{119}{15}$ & $\frac{120}{119}$ \\
\texttt{Sh\_GS?@\_??o??@\_???K????E?????K???} & 20 & 19 & $(1,1,1,1,1,1,1,1,1,1,2,2,2,2,2,2,2,2,3,9)$ & 8 & $\frac{1288}{165}$ & $\frac{165}{161}$ \\
\texttt{Sh\_K?E??K??@\_???o???@\_????K?????C} & 20 & 19 & $(1,1,1,1,1,1,1,1,1,2,2,2,2,2,2,2,2,2,2,9)$ & 9 & $\frac{179}{22}$ & $\frac{198}{179}$ \\
\texttt{Sh\_K?E??K??@\_???o???@\_????K??C???} & 20 & 19 & $(1,1,1,1,1,1,1,1,1,1,2,2,2,2,2,2,2,2,2,10)$ & 8 & $\frac{497}{66}$ & $\frac{528}{497}$ \\
\texttt{Sh\_K?E??K??@\_???o???@\_?A??C??C???} & 20 & 19 & $(1,1,1,1,1,1,1,1,1,1,1,2,2,2,2,2,2,2,2,11)$ & 7 & $\frac{539}{78}$ & $\frac{78}{77}$ \\
\texttt{SiPAAA??K??@\_???o???@\_????K?????C} & 20 & 19 & $(1,1,1,1,1,1,1,1,1,1,1,1,2,2,2,2,2,2,7,7)$ & 7 & $\frac{293}{42}$ & $\frac{294}{293}$ \\
\texttt{SiPAC?@\_??o??@\_???K????E?????K???} & 20 & 19 & $(1,1,1,1,1,1,1,1,1,1,1,1,2,2,2,2,2,2,6,8)$ & 7 & $\frac{2203}{315}$ & $\frac{2205}{2203}$ \\
\texttt{SiPC?E??K??@\_???o???@\_????K?????C} & 20 & 19 & $(1,1,1,1,1,1,1,1,1,1,1,2,2,2,2,2,2,2,5,8)$ & 8 & $\frac{491}{65}$ & $\frac{520}{491}$ \\
\texttt{SiPC?E??K??@\_???o???@\_????K??C???} & 20 & 19 & $(1,1,1,1,1,1,1,1,1,1,1,1,2,2,2,2,2,2,5,9)$ & 7 & $\frac{8047}{1155}$ & $\frac{8085}{8047}$ \\
\texttt{SiQ?GI??K??@\_???o???@\_????K?????C} & 20 & 19 & $(1,1,1,1,1,1,1,1,1,1,1,2,2,2,2,2,2,3,4,8)$ & 8 & $\frac{2557}{330}$ & $\frac{2640}{2557}$ \\
\texttt{SiQ?K?@\_??o??@\_???K????E?????K???} & 20 & 19 & $(1,1,1,1,1,1,1,1,1,1,1,2,2,2,2,2,2,2,4,9)$ & 8 & $\frac{16073}{2145}$ & $\frac{17160}{16073}$ \\
\texttt{SiQ?K?@\_??o??@\_???K????E??C??C???} & 20 & 19 & $(1,1,1,1,1,1,1,1,1,1,1,1,2,2,2,2,2,2,4,10)$ & 7 & $\frac{2652}{385}$ & $\frac{2695}{2652}$ \\
\texttt{Si\_GS?@?S??@\_???o???@\_????K?????C} & 20 & 19 & $(1,1,1,1,1,1,1,1,1,1,1,2,2,2,2,2,3,3,3,8)$ & 8 & $\frac{260}{33}$ & $\frac{66}{65}$ \\
\texttt{Si\_GS?@\_??o??@\_???K????E?????K???} & 20 & 19 & $(1,1,1,1,1,1,1,1,1,1,1,2,2,2,2,2,2,3,3,9)$ & 8 & $\frac{1258}{165}$ & $\frac{660}{629}$ \\
\texttt{Si\_K?E??K??@\_???o???@\_????K?????C} & 20 & 19 & $(1,1,1,1,1,1,1,1,1,1,2,2,2,2,2,2,2,2,3,9)$ & 9 & $\frac{175}{22}$ & $\frac{198}{175}$ \\
\texttt{Si\_K?E??K??@\_???o???@\_????K??C???} & 20 & 19 & $(1,1,1,1,1,1,1,1,1,1,1,2,2,2,2,2,2,2,3,10)$ & 8 & $\frac{6307}{858}$ & $\frac{6864}{6307}$ \\
\texttt{Si\_K?E??K??@\_???o???@\_?A??C??C???} & 20 & 19 & $(1,1,1,1,1,1,1,1,1,1,1,1,2,2,2,2,2,2,3,11)$ & 7 & $\frac{1838}{273}$ & $\frac{1911}{1838}$ \\
\texttt{SkE?K?@\_??o??@\_???K????E?????K???} & 20 & 19 & $(1,1,1,1,1,1,1,1,1,1,2,2,2,2,2,2,2,2,2,10)$ & 9 & $\frac{169}{22}$ & $\frac{198}{169}$ \\
\texttt{SkE?K?@\_??o??@\_???K????E??C??C???} & 20 & 19 & $(1,1,1,1,1,1,1,1,1,1,1,2,2,2,2,2,2,2,2,11)$ & 8 & $\frac{551}{78}$ & $\frac{624}{551}$ \\
\texttt{SkE?K?@\_??o??@\_???K??\_?A??C??C???} & 20 & 19 & $(1,1,1,1,1,1,1,1,1,1,1,1,2,2,2,2,2,2,2,12)$ & 7 & $\frac{251}{39}$ & $\frac{273}{251}$ \\
\texttt{SkE?K?@\_??o??@\_?O?C??\_?A??C??C???} & 20 & 19 & $(1,1,1,1,1,1,1,1,1,1,1,1,1,2,2,2,2,2,2,13)$ & 6 & $\frac{29}{5}$ & $\frac{30}{29}$ \\
\bottomrule
\end{tabular}
\end{table}

\section{観察}

\paragraph{反例の形}
145 個の反例のうち 131 個は木である、54 個は次数 $\ge n/2$ の頂点 (ハブ) をもつ、140 個は葉をもつ。
表 \ref{tab:ce} の次数列を見ると、反例はいずれも
「次数の大きい頂点が少数あり、残りはほとんど次数 2 以下」という
偏った次数分布をもつ。これは調和指数の定義から説明できる。
$H$ は各辺に $2/(d(u)+d(v))$ を配るので、次数の高い頂点に接する辺の
寄与は小さい。一方 $\mu^*$ は次数に鈍感で、長さ 2 の脚 (葉とその親) が
増えるたびに 1 ずつ増える。したがって「ハブ + 長さ 2 の脚」という形が
$\mu^*$ を大きく、$H$ を小さく保つ最も効率の良い構成になる。
実際、木の最小反例はこの形 (細分星) であり、これは Gupta \cite{gupta} の
指摘と一致する。

\paragraph{比の上界との隔たり}
観測された $\mu^*/H$ の最大値は $\frac{33}{28}$ である。
Gupta \cite{gupta} の木に対する上界 $\tfrac{3}{2}$、
一般グラフに対する $\mu^* > H/4$ という下界と比べると、
本稿の範囲の小さいグラフはいずれの極値からも遠い。
表 \ref{tab:fams} の最大比の列は位数とともに単調には動かず、
反例が存在しはじめる位数の近くで比が最大になっている。
上界に迫る族を小さい位数で見つけるのは難しく、
$n$ を増やしながら細分星の脚の本数を調整する構成的な族を
調べるほうが有望だと考えられる。

\paragraph{貪欲で足りる割合}
全 13,335,586 個のうち 84.7\% でしか厳密計算が必要にならなかった。
これは「NP 困難な量を含む不等式でも、片側であれば安価な発見的手法で
ほとんど片が付き、残りにだけ厳密解を使えばよい」という
証人つき検証の実務的な利点を示している。

\section{限界}

本稿の分類は $n \le 10$ の連結グラフと $n \le 20$ の木に限られる。
連結グラフについては McKay \cite{mckay} の完全リストが $n \le 10$ まで
しか公開されていないためであり、木については走査時間の都合である。
$n \ge 11$ の連結グラフに反例が何個あるかについて本稿は何も言わない。

また、反例でないグラフについて本稿が保証するのは
$\mu^*(G) \le H(G)$ という不等式だけで、$\mu^*(G)$ の値そのものは
(貪欲で片が付いたグラフでは) 確定していない。個々の $\mu^*$ をすべて
確定させようとすると全グラフで厳密計算が必要になり、証人つき検証の
利点は失われる。
ただし表 \ref{tab:fams} の最大比は例外で、厳密値として確定している。
証人が各グラフに $\mu^*(G) \le |M|$ という上界を与えるので
「族内のどのグラフも比 $|M|/H$ を超えない」ことが線形時間で確かめられ
(上から)、最大を達成すると主張されたグラフの $\mu^*$ だけを独立実装で
計算し直せば (下から)、両者が一致したときに最大値が確定する。
検証器はこの 2 つを各族で実際に行っている。


\section*{再現性}
本稿の主張はすべて、探索器とは独立に実装された検証器によって再検査されている。次のコマンドで再現できる。

\begin{center}\texttt{python -m mar verify p0003\_saturation\_harmonic}\end{center}

\begin{center}\begin{tabular}{ll}\toprule
証明書ダイジェスト & \texttt{3588557963d25d98} \\
生成日時 (UTC) & 2026-07-26T12:56:50+00:00 \\
Python & 3.10.11 \\
実行環境 & Windows 10 \\
リポジトリ版 & \texttt{00fa3fd} \\
乱数種 & 0 \\
探索時間 & 3139.3 秒 \\
\bottomrule\end{tabular}\end{center}


\begin{thebibliography}{99}
\bibitem{reverie} R. Davila ほか, In Reverie Together: Ten Years of Mathematical Discovery with a Machine Collaborator, arXiv:2507.17780, 2025. \url{https://arxiv.org/abs/2507.17780}
\bibitem{biyikoglu} T. Bıyıkoğlu, A note on a conjecture on the saturation number and the harmonic index, MATCH Commun. Math. Comput. Chem. 96 (2026) 1097--1099.
\bibitem{gupta} C. Gupta, Sharp bounds between the saturation number and the harmonic index, arXiv:2606.15761, 2026. \url{https://arxiv.org/abs/2606.15761}
\bibitem{fajtlowicz} S. Fajtlowicz, On conjectures of Graffiti II, Congr. Numer. 60 (1987) 187--197.
\bibitem{mckay} B. D. McKay, Combinatorial Data (graphs / trees). \url{https://users.cecs.anu.edu.au/~bdm/data/}
\bibitem{yannakakis} M. Yannakakis, F. Gavril, Edge dominating sets in graphs, SIAM J. Appl. Math. 38 (1980) 364--372.
\end{thebibliography}

\end{document}
